\begin{abstract}
\textbf{Context.} Learning-based vulnerability detectors are increasingly proposed as filters that decide which code a human need not inspect, yet they are evaluated by accuracy on random splits, which says little about the share of vulnerabilities such a filter silently lets through.
\textbf{Objective.} We study when an automated detector can be \emph{trusted to clear code}, by formulating vulnerability triage as a risk-control problem and testing its guarantees under the shifts that characterise software evolution: new projects, later vulnerabilities, different datasets and drift over time.
\textbf{Method.} We specialise conformal risk control to the clearance decision (finite-sample miss-rate guarantee, a high-probability variant, group-conditional calibration, a bound on the effect of distribution shift, an online variant with a deterministic guarantee under delayed feedback, certified flagging, and a cost-optimal budget) and evaluate it on three corpora (DiverseVul, Big-Vul, PrimeVul; 710\,000 deduplicated C/C++ functions) in 14 scenarios and three seeds, with four detectors, frozen and fine-tuned CodeBERT scorers and a rule-based analyser.
\textbf{Results.} Deduplication removes a quarter of Big-Vul; 43\% of Big-Vul and 77\% of PrimeVul also occur in DiverseVul, which inflates naive cross-dataset AUPRC by 3.5--5.6$\times$. Detection quality falls by 43--72\% outside the random split. On exchangeable data the guarantee holds (realised miss $9.3\%$ and $10.7\%$ for a $10\%$ budget); under shift it is violated by $2.7\times$ on average on DiverseVul and Big-Vul and by $2.0\times$ on PrimeVul (up to $37\%$), covariate-shift weighting repairs it only in some scenarios, and about 100 labelled target positives repair it, at a safe automation rate of 27--32\% instead of 54--55\% on DiverseVul and Big-Vul (26--51\% on PrimeVul). Adaptive conformal inference restores the long-run budget ($10.0$--$11.5\%$ versus $20$--$37\%$ for a static rule) from delayed outcome feedback alone, with a deterministic bound. Miss rates are very uneven across function sizes ($41\%$ vs.\ $2\%$) and CWEs, and 13.5\% of decisions flip under semantics-preserving rewrites. Frozen and fine-tuned CodeBERT scorers replicate the findings, and whitespace structure alone separates the classes with AUROC $0.89$ in Big-Vul ($0.84$ in PrimeVul), so raw-text models look valid and efficient until the dataset changes (miss rates up to $48\%$ at a $10\%$ budget). Delegation pays off only when a miss costs tens of reviews.
\textbf{Conclusion.} Reporting accuracy is not enough: trustworthy automation requires an explicit, monitored miss budget calibrated on the deployment distribution.
\begin{keyword}
vulnerability detection \sep conformal prediction \sep risk control \sep distribution shift \sep trustworthy automation \sep reproducibility
\end{keyword}
\end{abstract}
